\scanpage{229}
\mspage{103}
% archive sheet number: 228

De façon précise, \unsure{soit} $\pi$ un cardinal tel que
$\unsure{M'|N'}$ soit $\pi$-accessible et que $\varphi$ applique
$\unsure{\mathcal{M}|M}^\pi$ dans ${\mathcal{M}'}^\pi$ (\unsure{Par exemple|par exemple|Par Exemple}, pour cette dernière
condition, il suffit $\pi=2^c$ pour $c$ assez
grand, cf loc cit. 9.14 p. 153) Alors \deleted{$\varphi^{-1}(\unsure{N|M}$}
$M=\varphi^{-1}(M')$ est $\pi$-accessible.

Vérification immédiate.

On a donc le

\uline{Lemme} 8.2. Soit $A$ dans Cat. Alors $(A^\wedge,W_A)$
satisfait aux conditions préliminaires du
cor. \unsure{7|1}.4., ainsi qu'aux conditions suivantes

\[
 \deleted{\unsure{\mathcal C\subset|\mathrm{Le}}}
 \underset{\substack{\shortparallel\\\text{\unsure{ens.} des}\\\text{mono}\\\text{de } A^\wedge}}
 {\mathrm{Mon}(A^\wedge)}\subset\mathrm{Cof}_{W_A}
 \tag*{(8.2.1)}
\]

\[
 \deleted{\unsure{\mathrm i\in|\mathrm{ces}}}\quad \mathrm{Mon}\cap W_A
 \quad\text{stable\unsure{s} par coch. de base}
 \tag*{\unsure{(8.2.2)|(8.1.b)}}
\]

en d'autres termes (compte tenu que Mon est stable par coch. de base) pour
un carré cocartésien

\diagram{\begin{tikzcd}[ampersand replacement=\&,column sep=small]
 {(8.2.3)} \& A\arrow[rr,hook,"i"]\arrow[dr,"\alpha"'] \& \& B\arrow[dr,"\beta"] \& {}\\
 {} \& {} \& A'\arrow[rr,hook,"i'"] \& {} \& \deleted{\unsure{B|A}}B'
\end{tikzcd}}

avec $i$ mono, si $\alpha\in W_A$ (resp. $i\in W_A$),
on a $\beta\in W_A$ (resp. $i'\in W_A$). En fait
on a mieux : \deleted{le \unsure{dit|carré} carré est $W$-cocartésien}

\[
 \begin{aligned}
 &\deleted{i\in W_A\unsure{\Longleftrightarrow|\Longrightarrow}}\\
 &i\in W_A\unsure{\Longleftrightarrow|\Longrightarrow} i'\in W_A,\qquad
 \alpha\in W_A\Longleftrightarrow\beta\in W_A.
 \end{aligned}
 \tag*{(8.2.4)}
\]
